\documentclass[11pt]{article}
\usepackage[a4paper,margin=27mm]{geometry}
\usepackage{amsmath,amssymb,amsthm,hyperref}
\newcommand{\PSH}{\operatorname{PSH}}
\newcommand{\FS}{\operatorname{FS}}
\newcommand{\vol}{\operatorname{vol}}
\newcommand{\MA}{\operatorname{MA}}
\newtheorem{theorem}{Theorem}
\newtheorem{lemma}[theorem]{Lemma}
\title{An entire relative finite-energy equation for big exceptional classes}
\author{}
\date{6 October 2026}
\begin{document}
\maketitle
\begin{abstract}
For the nonnef big class $L=\pi^*A+cE$ on a blowup of a smooth
center, this paper constructs the whole genuine
relative finite-energy equation in every dimension. Exact section
division first identifies the entire genuine cone. A quadratic
ample-energy test excludes mass on the center. Actual moving model
intersections then define the original-carrier operator, with
weighted unbounded limits. Compact-coordinate envelopes and cutoff
comparisons prove both outer energy derivatives for every original
continuous direction. The intrinsic compensator dual energy yields
attained equations, comparison and uniqueness for all solutions.
The raw mass is $\operatorname{vol}(L)=A^n$. The sole equation
premise is the whole exact actual-ample theorem. The complete construction has passed independent review in this
precise signature. Original P16.5 for arbitrary big classes remains open;
zero-volume pseudoeffective classes are not asserted here.
\end{abstract}

\section{The exact domain and objects}
Let $k$ be any algebraically closed trivially valued field of
characteristic zero, $Y$ a smooth integral projective variety of
dimension $n\geq2$, and $Z\subset Y$ a nonempty smooth closed
center of codimension at least two. Let
$\pi:X=\operatorname{Bl}_Z Y\to Y$ and let $E$ be its
effective exceptional Cartier divisor. The source $X$ is smooth
integral projective. Fix an actual ample rational bundle $A$ on $Y$,
a rational $c>0$, and the actual big rational bundle
\[
 P=\pi^*A,\qquad L=P+cE.                                  \tag{1}
\]
The class is nonnef: a line in a geometric exceptional projective
space has $L$-degree $-c$. The center may have several disjoint
components; the single Cartier divisor $E$ and the same coefficient
$c$ are retained. No cardinality hypothesis on $k$ is imposed.

For an actual rational bundle $B$ a generically finite FS potential is
\[
 f=m^{-1}\max_j(\log|s_j|+\lambda_j),\qquad
 0\ne s_j\in H^0(mB),\quad \lambda_j\in\mathbf R.          \tag{2}
\]
Canonical trivial norms are used. The section system need not be
generated, and its poles are allowed. The genuine numerical cone is
exactly the class of one-net limits
\[
 \PSH(\theta)=\{u=\inf_i f_i:\ f_i\searrow\text{ on all points},
  \ f_i\in\FS^{\rm gf}(B_i),\ c_1(B_i)\to\theta,
                         u(x_0)>-\infty\}.               \tag{3}
\]
Here $x_0$ is the generic trivial valuation, and the index set may
have any cardinality. Normalization means $u(x_0)=0$. The topology
is the pointwise topology at divisorial valuations, not the all-point
topology. Birational proper $X,Y$ have the same set of these
valuations of the common function field.

The sole analytic premise $G_0$ is the exact whole actual-ample
theorem retained for original P16.1, applied separately to these
smooth projective $X$ and $Y$: finite continuous same-bundle FS
regularization of genuine ample potentials, decreasing closure with
finite generic endpoint, normalized divisorial compactness, finite
divisorial coordinates and all-point determination. No historical
general-big theorem or disputed abstract extended-real flattening
lemma is a premise. Necessary nef and intersection consequences are
proved below. Ordinary blowup, ample and normal-section facts are
also stated with their applications.

\section{Elementary genuine constructions}
An actual ample rational bundle has a generated power whose
weight-zero FS potential is identically zero. At every valuation-ring
center one generator is a unit and all generators have norm at most
one. Tensoring this system places a function (2) unchanged on the
sum of its bundle and that ample bundle. Hence every already
genuine function in a numerical class $\theta$ embeds, on a tail
of its defining net, in any fixed actual ample bundle $B$ with
$c_1(B)-\theta$ ample. Indeed $B-B_i$ is then actual ample.
This proves all-point determination in every genuine numerical class
by putting two functions in a common ample cone; it does not assert
that a new nonample limit is genuine.

Products of all section monomials after denominator clearing give
literal sums of FS potentials. A maximum of two FS potentials on
one actual bundle is FS there: take all required powers of each
system and then their weighted union. The same maximum is genuine
for two ample potentials, by taking maxima of their two finite
same-bundle regularizing nets, on the product-directed index set.
Consequently divisorial order for two already genuine potentials
in one numerical class implies all-point order: their maximum
in a common ample cone has the larger one's divisorial coordinates,
and $G_0$ gives equality everywhere.

We record the intersection construction needed to handle a limiting
class, with its pole issue resolved explicitly.
\begin{lemma}[Nef intersection]
For actual nef $Q$ and actual ample $H$ on either carrier,
\[
 \PSH(c_1(Q))=\bigcap_{\epsilon\in\mathbf Q_{>0}}
                         \PSH(Q+\epsilon H).             \tag{4}
\]
It is closed under arbitrary decreasing nets with finite limiting
generic value. Its normalized cone is compact in divisorial
coordinates and has one-net finite FS regularization in classes
tending to $c_1(Q)$.
\end{lemma}
\begin{proof}
The forward direction is the zero-system section embedding.
For the reverse direction let $u$ be in every displayed ample cone
and choose finite continuous $f_{\epsilon,j}\searrow u$ on
the actual bundle $Q+\epsilon H$. Consider all triples
$(\epsilon,j,a)$ with $a>0$, ordered by decreasing $\epsilon$,
decreasing $a$ and decreasing function $f_{\epsilon,j}+a$.
To dominate two triples choose smaller $\epsilon_3$ and smaller
positive $a_3$ than both shifts. The finite continuous function
$R=\min(f_{\epsilon_1,j_1}+a_1,f_{\epsilon_2,j_2}+a_2)$
strictly exceeds $u+a_3$. This is also strict at every pole of $u$,
since $R$ is finite. The open sets
\[
 \{\exp(f_{\epsilon_3,j}+a_3)-\exp R<0\}
\]
cover the compact analytic space. A finite subcover and directedness
give a common $j$ for which the comparison holds everywhere.
The triple system is therefore one decreasing FS net. Its classes
tend to $c_1(Q)$; at each point finite thresholds and arbitrarily
small shifts identify its infimum with $u$, also when $u=-\infty$.
This proves (4) in the exact definition (3).

An outer decreasing net belongs unchanged to every ample cone in
(4). Apply $G_0$ there and then (4). For compactness embed every
normalized ample cone $Q+\epsilon H$, $0<\epsilon\leq1$, into
the normalized cone $Q+H$; its compact image is closed in the
Hausdorff coordinate space. Their intersection is compact. All-point
determination in that common ample cone identifies the intersection
with the normalized original representatives in (4).
\end{proof}

For an actual ample $A$ on $Y$ the special intersection
$\bigcap_{\delta>0}\PSH((1+d\delta)A)=\PSH(A)$
is thus valid, using $Q=A$ and $H=dA$. Each member is understood
as the same actual all-point function, with the same normalization.

\section{Actual section division, including perturbations}
The blowup tautological bundle $\mathcal O_X(-E)$ is relatively
ample. Therefore, for a sufficiently large divisible integer $d$,
\[
 H=d\pi^*A-E\quad\hbox{is actual ample on }X.              \tag{5}
\]
Here this is the usual projective relative-ample plus pullback-ample
fact: embed a relatively very ample power into a projective bundle
over $Y$, tensor with a sufficiently large ample power on $Y$,
and use the absolute very ample embedding. Rational denominators
are cleared first.

For every integer $q\geq0$,
\[
                    \pi_*\mathcal O_X(qE)=\mathcal O_Y.  \tag{6}
\]
Indeed a section on $\pi^{-1}U$ is a rational function regular
off $Z\cap U$. It has no pole at any codimension-one point of
the normal variety $U$, since the removed center has codimension
at least two. A normal noetherian domain is the intersection of
its height-one valuation rings inside its fraction field, so it
extends to a regular function on $U$. Conversely its pullback is
regular and hence is a section of $\mathcal O_X(qE)$.
The identification as a section includes multiplication by the
canonical section $s_{qE}$ in local line-bundle frames. This proves
(6) as a sheaf assertion on each affine $U$, and they agree on
overlaps.

The projection formula, or the same local-frame proof, gives for
every actual rational bundle $B$ on $Y$ and rational $b\geq0$
the exact isomorphism
\[
 H^0(Y,mB)\longrightarrow H^0(X,m\pi^*B+mbE),\qquad
                  t\longmapsto s_{mbE}\pi^*t,             \tag{7}
\]
for all sufficiently divisible $m$. Every section on the right is
divisible by the prescribed exceptional factor; no numerical
equivalence is substituted for this actual identity.

For rational $0<\epsilon<c$ put
\[
 A_\epsilon=(1+d\epsilon)A,\quad
 N_\epsilon=(c-\epsilon)E,\quad
 L_\epsilon=L+\epsilon H=\pi^*A_\epsilon+N_\epsilon.
                                                               \tag{8}
\]
These are actual bundles and divisors, $A_\epsilon$ is ample,
$N_\epsilon$ is effective, and (7) applies to every actual section
of every sufficiently divisible power of $L_\epsilon$.
There is no nef or compactness assumption about $L_\epsilon$.
Let $\sigma_b=b\log|s_E|$; denominator clearing makes this an
actual one-section FS potential of $bE$ for rational $b\geq0$.
Its values are finite on divisorial points and zero at $x_0$.

\section{The entire cone theorem}
\begin{theorem}[Structural gate]
In the domain (1), the full genuine numerical cone is
\[
 \PSH(c_1(L))
       =\{\sigma_c+\pi^*w:w\in\PSH(Y,A)\}.              \tag{9}
\]
Every equality in (9) holds at all analytic points of the original
$X^{\rm an}$. On normalized cones this map is a homeomorphism
in divisorial coordinates. In particular the entire genuine big
cone is normalized compact and closed under arbitrary decreasing
nets with finite limiting generic value. Moreover
\[
 \PSH(c_1(\pi^*A))=\pi^*\PSH(Y,A),\qquad
                     \vol(L)=A^n.                         \tag{10}
\]
\end{theorem}
\begin{proof}
If $w\in\PSH(Y,A)$, its finite same-bundle regularizers pull
back to actual FS on $\pi^*A$, with the same decreasing
all-point limit. Tensoring the fixed $s_{mcE}$ gives one defining
FS net for $\sigma_c+\pi^*w$ on $L$ itself. Thus the right
side of (9) is genuine, with exactly the fixed exceptional ideal
powers rather than merely the same pole set.

Conversely let $u\in\PSH(c_1(L))$, normalized. Fix
$0<\epsilon<c$. If $f_i\in\FS^{\rm gf}(B_i)$ is its original
defining net, then $L_\epsilon-B_i$ is actual ample on a tail,
because its numerical class tends to $\epsilon c_1(H)$.
Tensor the zero FS systems to put the unchanged $f_i$ on
$L_\epsilon$. Clear the denominator of $N_\epsilon$ by
all products of each section system; this keeps the function
unchanged. Equation (7) divides every resulting section exactly
and gives actual FS functions $g_i$ on $Y,A_\epsilon$ with
\[
                 f_i=\sigma_{c-\epsilon}+\pi^*g_i        \tag{11}
\]
at all points. At poles this is simply the norm identity for a
product of sections; it uses no extended-real subtraction.

The common function-field divisorial evaluations of (11) are
finite, so $g_i$ decrease at every divisor. Their actual maximum
and $G_0$ on $Y,A_\epsilon$ upgrade this to all-point decrease.
The finite limiting generic value is zero. Thus $G_0$, or directly
the one FS net in (3), gives $w_\epsilon\in\PSH(Y,A_\epsilon)$
as their all-point infimum, normalized. Taking the literal infimum
in (11) yields
\[
 u=\sigma_{c-\epsilon}+\pi^*w_\epsilon
 \quad\hbox{on all of }X^{\rm an}.                         \tag{12}
\]
At points where the exceptional factor is infinite the equality
follows since every $f_i$ is $-\infty$; elsewhere its fixed finite
value commutes with the infimum. At a common divisor $v$ it reads
\[
 w_\epsilon(v)=u(v)-(c-\epsilon)\log|s_E|(v).              \tag{13}
\]

As rational $\epsilon\downarrow0$, the normalized $w_\epsilon$
all embed in a single normalized ample cone on $Y$. Compactness
gives a convergent subnet. For each rational $\delta>0$ the
subnet tail embeds in the compact closed normalized cone of
$(1+d\delta)A$, because $\epsilon<\delta$ eventually.
Its limit belongs to every such cone. The explicitly proved
intersection construction gives an actual $w\in\PSH(Y,A)$.
Equation (13) prescribes every finite divisorial limit as
$u(v)-c\log|s_E|(v)$. The already genuine function
$\sigma_c+\pi^*w$ therefore agrees with $u$ divisorially.
Both embed in one actual ample cone on $X$, where determination
gives equality at all points, including the poles. This proves (9).

On the common divisorial coordinates the map is fixed finite
translation, so it is a homeomorphism; injectivity follows by
ample determination on $Y$. Hence its entire normalized image is
compact. If $u_j=\sigma_c+\pi^*w_j$ decreases, (13) gives
divisorial decrease of $w_j$, and their ample maximum and
determination give all-point decrease. Their limiting generic
value is that of $u_j$. Apply $G_0$ to obtain
$w=\inf_j w_j$ on $Y$. Literal addition and pullback give
$\inf_j u_j=\sigma_c+\pi^*w$ at all points; the forward
construction supplies the required one actual FS defining net.

To prove the first part of (10), the nef-intersection lemma makes
every $v\in\PSH(c_1(\pi^*A))$ genuine. Its finite FS net from
that lemma, tensored with $s_{mcE}$, gives a genuine
$\sigma_c+v$ in class $L$. Equation (9) gives
$\sigma_c+v=\sigma_c+\pi^*w$ for some $w\in\PSH(Y,A)$.
Subtract the finite exceptional values on divisorial coordinates,
and use common-ample all-point determination on $X$ to get
$v=\pi^*w$. The reverse inclusion was already constructed by
pulling back the finite ample FS net. Finally (7) for $B=A,b=c$
gives $h^0(X,mL)=h^0(Y,mA)$ on all divisible powers. The ample
Hilbert polynomial has leading coefficient $A^n/n!$, proving
the raw volume in (10).
\end{proof}


\section{Main relative equation statement}
\begin{theorem}[Entire exceptional-class equation]
For every retained datum in (1), let $\sigma_c$ be the actual
canonical exceptional potential. The entire genuine cone is
$\sigma_c+\pi^*\PSH(Y,A)$, with intrinsic minimal reference
$\sigma_c$. Its relative finite-energy space is precisely
$\mathcal E_L^1=\sigma_c+\pi^*\mathcal E^1(Y,A)$.
The actual moving-intersection operator constructed below has raw
mass $V=A^n=\vol(L)>0$, independent of the regularizing net, and
satisfies mixed weighted decreasing convergence, comparison,
domination and uniqueness modulo constants. Its outer relative
energy has both directional derivatives in every
$f\in C^0(X^{\rm an})$, equal to $\int_X f\,\MA_L(u)$.

For every positive Radon measure $\nu$ on the original $X^{\rm an}$
with mass $V$, the equation $\MA_L(u)=\nu$ has an actual solution
$u\in\mathcal E_L^1$ if and only if the intrinsic compensator
dual energy in (24) is finite. Such a solution attains the relative
variational supremum, and all solutions differ by a constant.
\end{theorem}
\begin{proof}
The full structural argument precedes this statement. The following
sections construct the operator before the equation, prove the
polar exclusion and both outer derivatives, and finally establish
the exact target criterion, attainment and uniqueness.
\end{proof}

\section{The full ample equation premise, now used in its exact domain}
For the equation below we retain the whole mathematical body of $G_0$
on the actual ample $Y,A$, not just its structural part. In raw
normalization it gives energy $\mathsf E_A=V e_A$, $V=A^n>0$,
actual positive mixed measures of mass $V$ on $\mathcal E^1(Y,A)$,
decreasing unbounded and integrable weighted convergence, supporting
energy inequalities, comparison, domination and uniqueness. Every
positive Radon target of mass $V$ and finite intrinsic dual energy
has an actually attained energy maximizer and its ample equation.
The entire mixed operator is the genuine intersection operator.
In normalized conventions, $e_A(0)=0$, $e_A(w+a)=e_A(w)+a$,
and for $w(x_0)=0$ its $J$ is $-e_A(w)$, since the canonical
reference probability measure is $\delta_{x_0}$.
It also supplies
\[
 J(a w)\leq C_n a^2J(w),\quad0<a\leq1,\qquad
 \mathsf E_A^\downarrow(w+th)
  =\mathsf E_A(w)+t\int h\,\MA_A(w)+o(t)                 \tag{14}
\]
for every fixed $h\in C^0(Y^{\rm an})$, with both signs of $t$.
Its envelope theorem identifies this outer energy with energy of
the genuine envelope. Real PL functions are uniformly dense in
continuous functions and are differences of bounded genuine ample
potentials after increasing one actual ample bundle. These are
precise uses of the accepted actual-ample proof, Sections7--9.
They are not new nonample premises.

\section{Targets do not charge the blowup center}
\begin{lemma}
Every positive finite-dual-energy measure in the ample class $A$
gives zero mass to $Z^{\rm an}$. In particular this applies to
$\MA_A(w)$ for every $w\in\mathcal E^1(Y,A)$.
\end{lemma}
\begin{proof}
Work with its normalized probability $\mu/V$. Choose a sufficiently
large divisible $q$ and a nonzero section $h\in H^0(qA)$ vanishing
on $Z$. Such a section exists from the projective ample embedding,
or Serre generation of the coherent ideal of this proper center.
Put $\tau=q^{-1}\log|h|\leq0$, normalized at $x_0$.
For $T>0$, $b_T=\max(\tau,-T)$ is actual finite continuous FS on
$A$: tensor a generated zero FS system after clearing its power,
and take the weighted union. It is $-T$ on $Z^{\rm an}$,
has generic value zero, and $0\leq J(b_T)\leq T$.
For rational $0<a<1$, tensor the scaled system with the zero system
on $(1-a)A$; thus $ab_T$ is an actual finite-energy potential in
class $A$. If $m=(\mu/V)(Z^{\rm an})>0$, (14) gives
\[
 e_A(ab_T)-\int ab_T\,d(\mu/V)
                \geq (am-C_na^2)T.
\]
Choose $a$ small enough that the coefficient is positive. The dual
supremum would be infinite, a contradiction. Finite-energy outputs
have finite dual energy by the ample supporting inequality, so the
last assertion follows. This argument uses a single closed center
and arbitrary Radon measures, with no countable union of polar sets.
\end{proof}

Let $U_Y=Y^{\rm an}\setminus Z^{\rm an}$ and
$U_X=X^{\rm an}\setminus E^{\rm an}$. The algebraic blowup is
an isomorphism outside its center, so $\pi:U_X\to U_Y$ is an
analytic homeomorphism. Write $\iota:U_Y\to U_X$ for its inverse.
For a finite positive Radon $\mu$ with $\mu(Z^{\rm an})=0$ define
\[
             \mathcal L\mu=\iota_*(\mu|_{U_Y})             \tag{15}
\]
as a measure on the original compact $X^{\rm an}$, extended by zero
on $E^{\rm an}$. It is Radon: finite Radon restriction is tight on
the open $U_Y$, compact subsets push to compact subsets, and the
finite extension is regular. It has the same mass, gives zero mass
to $E^{\rm an}$, and $\pi_*\mathcal L\mu=\mu$.

The following continuity will be needed for actual unbounded limits.
If $\mu_j\to\mu$ weakly as positive Radon measures, with all their
center masses zero, including that of the limit, then
$\mathcal L\mu_j\to\mathcal L\mu$ weakly. For
$f\in C^0(X^{\rm an})$, $g=f\circ\iota$ is bounded continuous
on $U_Y$. Choose a continuous cutoff $0\leq\eta\leq1$ vanishing
on a neighborhood of $Z^{\rm an}$ with $\int(1-\eta)d\mu$
arbitrarily small. Extend $\eta g$ by zero; it is continuous on
$Y^{\rm an}$. Weak convergence applies to both $\eta g$ and
$1-\eta$. The bounded remainder is uniformly small eventually.
This proves the assertion for arbitrary nets. Applying it separately
to positive magnitude measures also proves the signed version when
those magnitudes converge weakly to a measure not charging the center.

\section{An intrinsic whole-cone relative energy and actual operator}
The structural theorem gives a unique $w\in\PSH(Y,A)$ for every
$u\in\PSH(c_1(L))$, with $u=\sigma_c+\pi^*w$.
Its canonical compensator is
\[
 \mathcal C(u)=\pi^*w\quad\hbox{on all of }X^{\rm an}.     \tag{16}
\]
Equivalently it is the unique genuine $P$-psh function whose sum
with $\sigma_c$ is $u$: divisorial subtraction is finite and
all-point determination proves uniqueness. This defines it even at
the common poles and makes no subtraction of two infinities.
The retained reference is the actual minimal potential $\sigma_c$;
normalized $u$ lies below it, since its normalized $w\leq0$.
Define on the entire genuine cone
\[
 \mathsf E_L(u)=\mathsf E_A(w),\qquad
 \mathcal E_L^1=\{u\in\PSH(c_1(L)):\mathsf E_L(u)>-\infty\}.
                                                               \tag{17}
\]
Thus the entire relative finite-energy domain is exactly
$\sigma_c+\pi^*\mathcal E^1(Y,A)$, with raw translation mass $V$.
Its energy is normalized to zero at $\sigma_c$, is monotone and
concave, and is continuous under every decreasing genuine net with
finite-energy endpoint. These assertions follow from the exact
ample properties by the bijection and all-point order already proved.
In particular this is a space of actual functions on the original
carrier, not a completion or an arbitrary choice on a modification.

We construct its operator before using any target equation. For a
finite nef PL/model potential $w$ on $Y$, the moving positive part
of $\sigma_c+\pi^*w$ is the pulled-back nef metric on $P$.
On a common model of $X$ and the model of $w$, projection of its
positive-part intersection cycle gives, component by component,
\[
        \MA_{L,\sigma_c}(\sigma_c+\pi^*w)
                       =\mathcal L\MA_A(w).               \tag{18}
\]
Here the genuine model formula is the sum of $b_F$ times the top
intersection of the moving nef part with the central components.
New exceptional components map into a smaller-dimensional central
subset, so the projection formula gives zero for their coefficient;
strict components give the original coefficient. The corresponding
divisorial analytic points are the same function-field valuations
and belong to $U_X,U_Y$, although their valuation-ring centers may
lie on the exceptional algebraic locus. The fixed exceptional factor is its total pullback on the common
model, including every vertical component; this entire fixed factor is
removed in the relative moving product.
Thus (18) is the actual intersection construction of this retained
reference product, not a prescribed inverse measure. Real model
weights are obtained by continuous finite intersection approximation.
The full ample proof supplies finite continuous FS regularizers
of each genuine finite-energy $w$, with finite PL/model refinements.
Their decreasing mixed intersection limits define the operator on
the whole domain (17). The center lemma and continuity of (15)
prove existence, independence of the regularizers and the formula
\[
                 \MA_L(u)=\mathcal L\MA_A(w),
                   \MA_L(u)(X^{\rm an})=V=\vol(L).         \tag{19}
\]
The identical model projection and limiting proof gives every mixed
relative product on this fixed class, with mass $V$. Each mixed output
has finite dual energy by the ample integrability criterion, so it
also charges no center. Weighted limits are genuine: for decreasing
finite-energy tests $z_j\to z$ and mixed ample outputs
$\mu_j\to\mu$, subtract one common upper constant on a tail.
The tests still decrease and their negative weighted measures are
positive. The ample theorem gives their weak convergence; the limit
charges no center because $\mu$ does not and $z$ is integrable.
Apply the signed continuity statement after (15), and restore constants
by the weighted-measure identity. Generic constants converge, but
normalizing each stage by its own generic value is not used to claim
decrease. This proves
weighted mixed convergence for the actual compensators
$\mathcal C(v_j)=\pi^*z_j$, on $X$. Integrals of bare $u_j$
include the fixed logarithmic factor and need not be finite; all
relative weights here are explicitly the compensators in (16).

\section{Genuine envelopes and both continuous-direction derivatives}
We first explain a useful order fact on either carrier. If $a,b$
are genuine potentials and $f$ is continuous, the constraint
$a\leq b+f$ on all divisorial points implies it on all analytic
points. Approximate $f$ uniformly by real PL functions. Express
each as a difference of bounded genuine ample potentials. Tensor
and put the two compared sums in a common sufficiently large
actual ample cone; divisorial order and $G_0$ give all-point order.
The uniform errors then tend to zero. The same proof treats a
continuous obstacle without $b$, comparing the two sides in a
common ample cone even if their numerical classes initially differ.
Every step uses already genuine sums with one product-directed
defining net.

Fix $u\in\mathcal E_L^1$ and $f\in C^0(X^{\rm an})$, and write
$M=\|f\|_\infty$. The admissible set below $u+tf$ is nonempty,
since it contains $u-|t|M$. Replace each competitor by its maximum
with that lower bound. Its generic value lies in one fixed compact
interval, and its normalized cone is compact by the structural
theorem. The constraints are closed divisorial-coordinate
inequalities by the preceding order fact. Finite maxima stay
admissible. The net of finite maxima has a cluster point in this
compact set, and its coordinates are the supremum of the
competitor coordinates. All-point order makes that point the
greatest admissible genuine function. Denote it
$P_L(u+tf)$. It is between $u-|t|M$ and $u+|t|M$, so it is
in the entire domain (17).

For an upper bounded usc obstacle $h$ with a genuine lower anchor,
define the outer extension
of the competitor energy by infimum over its continuous majorants.
In particular
\[
 \mathsf E_L^\downarrow(u+tf)=\mathsf E_L(P_L(u+tf)).       \tag{20}
\]
To verify this rather than postulate it, use the usual decreasing
net of continuous majorants of $h$ on a compact Hausdorff space.
Restrict to a fixed first majorant. Their greatest genuine minorants
exist by the same compact-coordinate maximum argument, stay above
$u-|t|M$ when $h=u+tf$, and decrease to the genuine greatest
minorant of $h$ by the proved decreasing closure. The corresponding
ample remainders decrease too. Ample decreasing energy continuity
gives (20). The order comparison with a continuous obstacle was
established above, so no nonample envelope theorem is hidden here.

Now let $u=\sigma_c+\pi^*w$ and $\mu=\MA_A(w)$. Put
$g=f\circ\iota$ on $U_Y$. For a cutoff $\eta$ as after (15)
define continuous functions on $Y^{\rm an}$
\[
 g_\pm=\eta g\ \pm M(1-\eta),\qquad
          g_-\leq g\leq g_+\quad\hbox{on }U_Y.            \tag{21}
\]
For $t>0$, genuine ample envelopes and divisorial/all-point order
give
\[
 \mathsf E_A(P_A(w+t g_-))
 \leq\mathsf E_L(P_L(u+tf))
 \leq\mathsf E_A(P_A(w+t g_+)).                            \tag{22}
\]
The lower envelope supplies an admissible potential on $X$;
at $E^{\rm an}$ its inequality is automatic because the fixed
factor is $-\infty$. For the upper envelope, an admissible
remainder satisfies its inequality on $U_Y$, hence on every
divisor; the order fact upgrades it to the all-point ample
obstacle constraint. For $t<0$ interchange $g_-$ and $g_+$.
Apply the two-sided ample expansion (14) at fixed cutoff, then
let the center-neighborhood measure tend to zero. The difference
of the derivative bounds is at most $2M\int(1-\eta)d\mu$.
Consequently
\[
 \left.\frac{d}{dt}\right|_{t=0^+}\mathsf E_L^\downarrow(u+tf)
 =\left.\frac{d}{dt}\right|_{t=0^-}\mathsf E_L^\downarrow(u+tf)
 =\int_{X^{\rm an}}f\,\MA_L(u).                          \tag{23}
\]
Time tends to zero before the cutoff changes. This proves every
original continuous direction; directions supported near the
exceptional locus were not discarded by a descent assumption.

\section{Intrinsic targets, attained equations and all solutions}
For a positive Radon $\nu$ of mass $V$ on the original $X^{\rm an}$,
define its intrinsic relative dual energy by
\[
 \mathsf E_L^\vee(\nu)=
 \sup_{u\in\mathcal E_L^1}
       \left\{\mathsf E_L(u)-\int_X\mathcal C(u)\,d\nu\right\}.
                                                               \tag{24}
\]
The integral is the canonical compensator (16), including its
specified values at the common poles; it is not an undefined bare
difference. Constants cancel exactly because the target mass is $V$.
Writing $\mu=\pi_*\nu$, (16)--(17) give the exact identity
\[
                        \mathsf E_L^\vee(\nu)
                              =\mathsf E_A^\vee(\mu).     \tag{25}
\]
Thus finiteness is equivalent to integrability of every actual
relative finite-energy compensator, by the full ample intrinsic
criterion. The center lemma forces $\mu(Z^{\rm an})=0$,
so necessarily $\nu(E^{\rm an})=0$ and $\nu=\mathcal L\mu$.
This exclusion follows from (24); it is not an independent support
restriction substituted for the target condition.

If (24) is finite, the exact ample variational theorem gives an
actual $w\in\mathcal E^1(Y,A)$ maximizing
$\mathsf E_A(w)-\int w\,d\mu$ and satisfying $\MA_A(w)=\mu$.
Its one-net pullback with the fixed actual exceptional sections is
an actual $u=\sigma_c+\pi^*w\in\mathcal E_L^1$ on $X$.
Equations (19) and (25) give
\[
                              \MA_L(u)=\nu.              \tag{26}
\]
The objective is identically the ample objective, so this is an
actually attained relative maximizer as well. Alternatively (23)
gives the two variational signs at that maximizer and identifies
the equation on all continuous tests. Conversely any solution
has finite dual energy by (25) and the ample supporting inequality.
The operator was constructed by moving intersections before this
argument, so (26) does not define a measure through an inverse.

For any two actual finite-energy functions $u,v$, their raw
comparison inequality and strict-winning locality follow from
the ample inequalities on $U_Y$ via the homeomorphism and (19).
Both measures vanish on $E^{\rm an}$, where both functions have
the fixed pole. If $u\leq v$ almost everywhere for $\MA_L(v)$,
subtract only the finite fixed factor on $U_X$; ample domination
gives their remainder order everywhere on $Y$, and hence $u\leq v$
on all of $X$. If $\MA_L(u)=\MA_L(v)$, push forward and use
ample uniqueness to obtain $u-v$ equal to one additive constant
in the metric sense $u=v+C$. This formulation handles shared poles
without undefined differences. It proves uniqueness for all actual
solutions in the entire domain (17), not just the constructed one.

\section{Exact scope of the equation theorem}
The argument supplies the entire genuine cone, its normalized
compactness and arbitrary-net closure, the full relative energy,
raw moving-intersection operator and weighted limits, both energy
derivatives in every original continuous direction, the intrinsic
target criterion, comparison, actual attainment and all-solution
uniqueness for the nonnef big blowup family (1). All field cardinalities
and all dimensions $n\geq2$ are retained. Its mass is
$\vol(L)=A^n$, which need not equal $L^n$.
The exact whole actual-ample $G_0$ proof on $Y,A$ is the only
equation grant; the structural proof on $X$ is included here in full.
The complete structural and equation proof has passed independent
review, relative to that exact actual-ample premise. It contributes
one accepted supporting paper; it closes no additional original atom.

Original P16.5 also requires arbitrary big rational bundles and the
additional pseudoeffective zero-volume range. Not every such bundle
has presentation (1). Neither those full classes nor a probability
theory at zero volume are asserted. The full unchanged original
P16.5 remains open. The frozen candidate and canonical independent review are retained
separately. None of the new editorial clarifications changes an active
research input.

\end{document}
